% !TEX program = xelatex
% =====================================================================
% LLM Theory Seminar -- English Study Notes
% Compile with:
%   xelatex FILENAME.tex
%   xelatex FILENAME.tex
% or: latexmk -xelatex FILENAME.tex
% =====================================================================
\documentclass[11pt,a4paper,oneside,openany]{memoir}
\usepackage{amsmath,amssymb,amsthm,mathtools,bm}
\usepackage{booktabs,longtable,array,tabularx,makecell}
\usepackage{enumitem}
\usepackage[most]{tcolorbox}
\usepackage[dvipsnames]{xcolor}
\usepackage[unicode,bookmarks=false]{hyperref}
\usepackage{cleveref}
\usepackage{needspace}
\setlrmarginsandblock{27mm}{27mm}{*}
\setulmarginsandblock{24mm}{28mm}{*}
\checkandfixthelayout
\setlength{\parindent}{0pt}
\setlength{\parskip}{0.48em}
\setlength{\emergencystretch}{3em}
\hbadness=10000
\setlength{\headheight}{15pt}
\linespread{1.055}\selectfont
\setlist[itemize]{topsep=0.28em,itemsep=0.12em,parsep=0pt,leftmargin=1.55em}
\setlist[enumerate]{topsep=0.28em,itemsep=0.16em,parsep=0pt,leftmargin=1.75em}
\definecolor{NoteBlue}{HTML}{294E6B}
\definecolor{NoteBlueLight}{HTML}{F2F6F9}
\definecolor{NoteGray}{HTML}{5A6268}
\definecolor{NoteGrayLight}{HTML}{F6F6F4}
\definecolor{NoteWarm}{HTML}{7A6545}
\definecolor{NoteWarmLight}{HTML}{FAF7F0}
\hypersetup{colorlinks=true,linkcolor=NoteBlue,citecolor=NoteBlue,urlcolor=NoteBlue}
\makepagestyle{seminar}
\makeoddhead{seminar}{\footnotesize LLM Theory Seminar}{}{\footnotesize NTK, Lazy Training \& Feature Learning}
\makeoddfoot{seminar}{}{\footnotesize\thepage}{}
\makeheadrule{seminar}{\textwidth}{0.35pt}
\pagestyle{seminar}
\aliaspagestyle{chapter}{seminar}
\makechapterstyle{seminarnotes}{
  \setlength{\beforechapskip}{8pt}
  \setlength{\midchapskip}{8pt}
  \setlength{\afterchapskip}{22pt}
  \renewcommand{\chapterheadstart}{\vspace*{\beforechapskip}}
  \renewcommand{\chapnamefont}{\normalfont\small\bfseries}
  \renewcommand{\chapnumfont}{\normalfont\bfseries\fontsize{34}{38}\selectfont\color{NoteBlue}}
  \renewcommand{\chaptitlefont}{\normalfont\huge\bfseries\color{black}}
  \renewcommand{\printchaptername}{}
  \renewcommand{\chapternamenum}{}
  \renewcommand{\printchapternum}{\chapnumfont\thechapter}
  \renewcommand{\afterchapternum}{\par\nobreak\color{black}\vskip\midchapskip}
  \renewcommand{\printchaptertitle}[1]{\chaptitlefont ##1}
  \renewcommand{\afterchaptertitle}{\par\nobreak\vskip\afterchapskip}
}
\chapterstyle{seminarnotes}
\setsecheadstyle{\Large\bfseries}
\setsubsecheadstyle{\large\bfseries}
\setsubsubsecheadstyle{\normalsize\bfseries}
\setbeforesecskip{-1.3\baselineskip plus -0.2\baselineskip minus -0.1\baselineskip}
\setaftersecskip{0.55\baselineskip}
\setbeforesubsecskip{-0.9\baselineskip plus -0.15\baselineskip minus -0.1\baselineskip}
\setaftersubsecskip{0.35\baselineskip}
\newcommand{\R}{\mathbb{R}}
\newcommand{\E}{\mathbb{E}}
\newcommand{\norm}[1]{\left\lVert #1\right\rVert}
\newcommand{\inner}[2]{\left\langle #1,#2\right\rangle}
\newcommand{\grad}{\nabla}
\DeclareMathOperator*{\argmin}{arg\,min}
\DeclareMathOperator*{\argmax}{arg\,max}
\DeclareMathOperator{\rank}{rank}
\DeclareMathOperator{\Tr}{Tr}
\tcbset{seminarbox/.style={enhanced,breakable,sharp corners,boxrule=0pt,frame hidden,left=3.2mm,right=3mm,top=1.8mm,bottom=1.8mm,before={\par\Needspace{5\baselineskip}},before skip=0.8em,after skip=0.8em,fonttitle=\bfseries\small,coltitle=black,attach title to upper={\quad}}}
\newtcolorbox{keybox}[1][]{seminarbox,colback=NoteBlueLight,borderline west={1.8pt}{0pt}{NoteBlue},title={Key point#1}}
\newtcolorbox{paperbox}[1][]{seminarbox,colback=NoteGrayLight,borderline west={1.8pt}{0pt}{NoteGray},title={From the original paper#1}}
\newtcolorbox{suppbox}[1][]{seminarbox,colback=NoteWarmLight,borderline west={1.8pt}{0pt}{NoteWarm},title={Supplementary derivation#1}}
\newtcolorbox{warningbox}[1][]{seminarbox,colback=NoteGrayLight,borderline west={1.8pt}{0pt}{NoteGray},title={Caution#1}}
\newtcolorbox{presentbox}{seminarbox,colback=white,borderline west={2.2pt}{0pt}{black!70},fontupper=\footnotesize,top=1.2mm,bottom=1.2mm,before={\par},before skip=0.5em,after skip=0.5em,title={How to say it in the presentation}}
\theoremstyle{definition}
\newtheorem{definition}{Definition}[chapter]
\newtheorem{proposition}{Proposition}[chapter]
\newtheorem{theorem}{Theorem}[chapter]
\begin{document}
\begin{titlingpage}
\thispagestyle{empty}
\vspace*{1.2cm}
{\small\bfseries LLM Theory Seminar \hfill Week 3\par}
\vspace{1.4cm}
{\HUGE\bfseries Neural Tangent Kernel, Lazy Training,\\[0.18em]and Feature Learning\par}
\vspace{0.55cm}
{\large English seminar study notes\par}
\vspace{0.9cm}
\rule{\textwidth}{0.7pt}
\vspace{1.1cm}
\begin{minipage}{0.86\textwidth}
\small
These notes are an English companion to the seminar handout.  The emphasis is on
mathematical derivations that can be followed line by line, on separating theorem
statements from interpretation, and on identifying the assumptions under which each
claim is valid.
\end{minipage}
\vfill
{\small\textbf{Primary readings}\\[0.25em]
Jacot, Gabriel, Hongler, \textbf{Neural Tangent Kernel}, NeurIPS 2018.\\Lee et al., \textbf{Wide Neural Networks of Any Depth Evolve as Linear Models Under Gradient Descent}, NeurIPS 2019.\\Chizat, Oyallon, Bach, \textbf{On Lazy Training in Differentiable Programming}, NeurIPS 2019.
\par}
\vspace{0.8cm}
{\small September 2026\par}
\end{titlingpage}
\tableofcontents*
\clearpage

\chapter{Notation and the central question}
\section{Notation map}
Let $f_\theta(x)$ be a neural network, $\theta\in\R^p$ its parameters, and $X=(x_1,\ldots,x_n)$ the training inputs.  Write
\[
f_\theta(X)=\begin{bmatrix}f_\theta(x_1)&\cdots&f_\theta(x_n)\end{bmatrix}^\top,
\qquad
J_\theta(X)=\frac{\partial f_\theta(X)}{\partial \theta}.
\]
The empirical neural tangent kernel is
\[
K_\theta(X,X)=J_\theta(X)J_\theta(X)^\top.
\]
For square loss we use the residual $r_\theta=f_\theta(X)-y$.

\begin{center}\small
\begin{tabularx}{\textwidth}{>{\bfseries\raggedright\arraybackslash}p{0.22\textwidth}XXX}
\toprule
Concept & Jacot et al. & Lee et al. & Chizat et al.\\
\midrule
network & $f_\theta$ & $f(\theta,x)$ & $h(w)$ / scaled model\\
tangent kernel & $\Theta_\infty$ & $\hat\Theta_t$ & tangent model\\
linearization & implicit in kernel limit & explicit Taylor model & lazy regime\\
feature learning & absent in strict NTK limit & deviation from linearization & non-lazy regime\\
\bottomrule
\end{tabularx}
\end{center}

\section{The question of the week}
Mean-field theory showed that infinite width can preserve nontrivial feature motion.  NTK theory asks about a different limit: can a wide network remain so close to initialization that training is described by its first-order Taylor expansion?

The crucial distinction is
\[
\text{width }m\to\infty
\quad\not\Rightarrow\quad
\text{a unique limiting dynamics}.
\]
The parameterization, output scaling, learning rate, and initialization jointly decide whether the limit is kernel-like or feature-learning.

\begin{presentbox}
The NTK is not introduced as an arbitrary similarity measure.  It appears exactly when parameter gradient flow is pushed through the network Jacobian into function space.
\end{presentbox}

\chapter{From parameter gradient flow to the NTK}
\section{Function-space dynamics}
For square loss
\[
\mathcal L(\theta)=\frac12\|f_\theta(X)-y\|^2=\frac12\|r_\theta\|^2,
\]
parameter gradient flow is
\[
\dot\theta_t=-\nabla_\theta\mathcal L(\theta_t).
\]
By the chain rule,
\[
\nabla_\theta\mathcal L(\theta_t)=J_t^\top r_t.
\]
Therefore
\[
\dot\theta_t=-J_t^\top r_t.
\]
Now differentiate the predictions:
\begin{align*}
\frac{d}{dt}f_{\theta_t}(X)
&=J_t\dot\theta_t\\
&=-J_tJ_t^\top r_t\\
&=-K_t r_t.
\end{align*}
Since $y$ is fixed,
\[
\boxed{\dot r_t=-K_t r_t.}
\]
This identity is exact for any differentiable network under parameter gradient flow.

\section{If the kernel is constant}
If $K_t\equiv K_0$, then
\[
\dot r_t=-K_0r_t,
\]
so
\[
\boxed{r_t=e^{-K_0t}r_0.}
\]
If $K_0=Q\Lambda Q^\top$, each eigenmode decays as $e^{-\lambda_i t}$.  Thus the kernel spectrum determines the optimization time scales.

\section{General losses}
For a loss $\ell(f(X),y)$,
\[
\dot f_t(X)=-K_t\nabla_f\ell(f_t(X),y).
\]
The NTK therefore acts as a data-dependent preconditioner in function space.

\begin{presentbox}
First derive $\dot f=-K_t\nabla_f\ell$ before discussing infinite width.  The infinite-width theorem is then simply the claim that, under a particular scaling, $K_t$ converges to a deterministic kernel and barely changes during training.
\end{presentbox}

\chapter{Jacot et al.: the infinite-width NTK}
\section{Kernel at initialization}
For fully connected networks with standard independent initialization, layerwise covariances converge to deterministic recursions.  Schematically,
\[
\Sigma^{(0)}(x,x')=\frac1d x^\top x',
\]
and if $(u,v)$ is jointly Gaussian with covariance determined by $\Sigma^{(\ell)}$,
\[
\Sigma^{(\ell+1)}(x,x')
=\E[\sigma(u)\sigma(v)].
\]
The derivative covariance is
\[
\dot\Sigma^{(\ell+1)}(x,x')
=\E[\sigma'(u)\sigma'(v)].
\]
For a standard fully connected parameterization, the limiting NTK obeys a recursion of the form
\[
\Theta^{(\ell+1)}
=\Sigma^{(\ell+1)}+\Theta^{(\ell)}\dot\Sigma^{(\ell+1)}.
\]
The exact base case and scaling depend on the architecture convention.

\section{Kernel constancy during training}
In the NTK parameterization, sufficiently wide networks move only a small distance in parameter space over $O(1)$ training time.  As width tends to infinity,
\[
K_t\longrightarrow \Theta_\infty
\]
uniformly over finite time intervals under the theorem assumptions.  The training dynamics then becomes
\[
\dot f_t(X)=-\Theta_\infty\nabla_f\ell(f_t(X),y).
\]
For square loss this is a linear ODE.

\begin{warningbox}
``The NTK is constant'' is an infinite-width statement under a specific scaling.  Finite networks generally have $K_t\neq K_0$, and the amount of kernel evolution is itself a useful diagnostic of feature learning.
\end{warningbox}

\begin{presentbox}
Jacot et al. identify a regime in which a nonlinear neural network has a nonlinear forward map but nearly linear training dynamics around initialization.  The nonlinearity survives in the deterministic kernel $\Theta_\infty$; the learned features do not substantially evolve.
\end{presentbox}

\chapter{Linearization and kernel regression}
\section{First-order Taylor model}
Linearize around initialization $\theta_0$:
\[
f_{\mathrm{lin}}(x;\theta)
=f_{\theta_0}(x)+J_0(x)(\theta-\theta_0).
\]
Training this linear model by gradient flow yields the same function-space equation with fixed kernel $K_0$.

\section{Minimum-norm solution}
Suppose the linearized model interpolates the training set.  Let $\Delta\theta=\theta-\theta_0$ and solve
\[
J_0\Delta\theta=y-f_0(X).
\]
Gradient flow from $\Delta\theta=0$ selects the minimum Euclidean-norm solution
\[
\Delta\theta_\star
=J_0^\top(J_0J_0^\top)^+\bigl(y-f_0(X)\bigr).
\]
For a test point $x$,
\begin{align*}
f_\star(x)
&=f_0(x)+J_0(x)\Delta\theta_\star\\
&=f_0(x)+K_0(x,X)K_0(X,X)^+\bigl(y-f_0(X)\bigr).
\end{align*}
This is the ridgeless kernel-regression predictor with the empirical NTK.

\section{Lee et al.: wide networks follow their linearization}
Lee et al. make the linearization statement explicit: under suitable parameterization, width, learning-rate, and time assumptions, the trained wide network remains close to the output of its first-order Taylor model.  One consequence is that optimization and prediction can be approximated by a kernel method without solving the nonlinear parameter dynamics.

\begin{presentbox}
The bridge from NTK dynamics to kernel regression is not a metaphor.  The first-order Taylor model has features $J_0(x)$, and its Gram matrix is exactly $J_0J_0^\top=K_0$.
\end{presentbox}

\chapter{What ``lazy training'' means}
\section{Small parameter displacement}
A differentiable model $h(w)$ is lazy when training stays in a neighborhood where
\[
h(w)\approx h(w_0)+Dh(w_0)(w-w_0)
\]
is accurate over the relevant trajectory.  A useful dimensionless indicator discussed by Chizat--Oyallon--Bach compares initial error, second derivatives, and first derivatives, schematically
\[
\kappa_h(w_0)
\sim
\|h(w_0)-y^\star\|
\frac{\|D^2h(w_0)\|}{\|Dh(w_0)\|^2}.
\]
A small value means that the model can fit the target using a displacement small enough that curvature of the parameter-to-output map remains negligible.

\section{Scaling can force laziness}
Consider a scaled model whose output sensitivity grows with a scale parameter $\alpha$.  Then an $O(1)$ change in output may require only an $O(1/\alpha)$ parameter displacement.  As $\alpha\to\infty$, the nonlinear model can therefore be forced into its tangent approximation even if the architecture itself is unchanged.

\section{Lazy does not mean slow}
The adjective refers to feature motion, not necessarily to optimization speed.  Kernel eigenvalues can be large, so a lazy model may fit rapidly while still learning almost no new representation.

\begin{presentbox}
Lazy training is best defined geometrically: the trajectory is short enough that the tangent features barely change.  It is not synonymous with ``wide,'' ``small learning rate,'' or ``slow convergence'' by themselves.
\end{presentbox}

\chapter{Feature learning and the scaling dichotomy}
\section{Mean-field versus NTK scaling}
For a schematic two-layer network
\[
f_m(x)=\alpha_m\sum_{i=1}^m a_i\sigma(w_i^\top x),
\]
choices such as $\alpha_m\propto m^{-1/2}$ and $\alpha_m\propto m^{-1}$ lead to different balances between output magnitude and parameter motion.  In the kernel-like scaling, a small displacement of many parameters produces an $O(1)$ output change.  In the mean-field scaling, neurons can move $O(1)$ and the feature distribution evolves.

\section{Yang--Hu: feature learning in infinite-width limits}
Tensor-program analyses make the scaling question systematic.  Under broad architecture classes and stable parameterizations, infinite-width limits can fall into kernel-like regimes or genuine feature-learning regimes.  The point is not that every parameterization is exactly one of two textbook formulas, but that nontrivial feature learning requires balancing layerwise updates so internal representations change at $O(1)$ scale.

\section{$\mu$P as a practical parameterization principle}
Maximal Update Parameterization ($\mu$P) chooses width-dependent parameter and learning-rate scalings so that updates remain nontrivial and stable as width changes.  It is designed to preserve feature-learning behavior rather than collapse to a fixed kernel.

\begin{warningbox}
Do not read the kernel/feature-learning distinction as a statement that feature learning is always superior.  Kernel limits are analytically clean and can perform very well.  The theoretical distinction concerns what the representation is allowed to change, not a universal ranking of predictors.
\end{warningbox}

\begin{presentbox}
The important axis is not finite width versus infinite width.  It is frozen tangent features versus evolving features.  Both behaviors can survive a width limit, depending on the scaling.
\end{presentbox}

\chapter{How to compare the regimes}
\section{A compact table}
\begin{center}\small
\begin{tabularx}{\textwidth}{>{\bfseries\raggedright\arraybackslash}p{0.22\textwidth}XXX}
\toprule
 & NTK / lazy & Mean-field & Finite feature learning\\
\midrule
parameter motion & vanishing/small & $O(1)$ after rescaling & architecture dependent\\
kernel evolution & vanishes in limit & generally nontrivial & measurable\\
representation & effectively fixed & distribution evolves & evolves\\
limit equation & kernel ODE & transport PDE & nonlinear SGD/GD\\
analysis tool & RKHS, spectra & Wasserstein, PDE & nonlinear optimization\\
\bottomrule
\end{tabularx}
\end{center}

\section{Diagnostics in experiments}
Common empirical questions include:
\begin{itemize}
\item How much does $\|K_t-K_0\|$ change?
\item How far do parameters move relative to initialization?
\item Do hidden representations align with new task structure?
\item Does a frozen-feature or kernel baseline match the trained network?
\end{itemize}
No single diagnostic is a universal definition, but together they distinguish the regimes.

\section{What NTK theory does not say}
NTK convergence does not prove that SGD finds the same solution at realistic finite width, that feature learning is absent in practical networks, or that a kernel predictor has the same out-of-distribution behavior as the finite network.

\begin{presentbox}
The cleanest seminar message is: NTK theory is a controlled linearization theorem.  It becomes powerful precisely because it turns training into a fixed-kernel dynamical system; its limitation is the same simplification---the learned representation is frozen to first order.
\end{presentbox}

\chapter{Presentation checklist}
\section{Eight equations to keep}
\begin{align*}
&\dot\theta=-J^\top r,\\
&K_\theta=JJ^\top,\\
&\dot r=-K_t r,\\
&r_t=e^{-K_0t}r_0\quad(K_t=K_0),\\
&f_{\mathrm{lin}}(x)=f_0(x)+J_0(x)(\theta-\theta_0),\\
&f_\star(x)=f_0(x)+K_0(x,X)K_0(X,X)^+(y-f_0(X)),\\
&\Theta^{(\ell+1)}=\Sigma^{(\ell+1)}+\Theta^{(\ell)}\dot\Sigma^{(\ell+1)},\\
&\rho_t\text{ evolves nontrivially in the mean-field regime.}
\end{align*}

\section{Common mistakes}
\begin{itemize}
\item Saying that the NTK identity $K=JJ^\top$ is itself an infinite-width result; it is finite-width and exact.
\item Saying the kernel is always constant during neural-network training.
\item Equating infinite width with a unique kernel limit.
\item Confusing small parameter displacement with small function change.
\item Treating a representational theorem as a claim about what SGD necessarily learns.
\end{itemize}

\begin{presentbox}
Thirty-second conclusion: pushing parameter gradient flow through the Jacobian produces the NTK exactly.  In one infinite-width scaling the kernel freezes and the network follows its linearization; in other scalings, including mean-field and $\mu$P-style regimes, features continue to move at leading order.
\end{presentbox}

\chapter*{Bibliography}
\addcontentsline{toc}{chapter}{Bibliography}
\begin{itemize}
\item Jacot, A., Gabriel, F., and Hongler, C. \emph{Neural Tangent Kernel: Convergence and Generalization in Neural Networks}. NeurIPS, 2018.
\item Lee, J. et al. \emph{Wide Neural Networks of Any Depth Evolve as Linear Models Under Gradient Descent}. NeurIPS, 2019.
\item Chizat, L., Oyallon, E., and Bach, F. \emph{On Lazy Training in Differentiable Programming}. NeurIPS, 2019.
\item Yang, G. and Hu, E. J. \emph{Tensor Programs IV: Feature Learning in Infinite-Width Neural Networks}. ICML, 2021.
\end{itemize}
\end{document}
